\section{Training objectives}
\label{sec:hamiltonian-loss}

\subsection{Base supervised loss}

Let
$q^\star=\G^\star(\eta;h)\xi$ denote the stored source-protocol reference
target defined in \cref{sec:data}, and let
$q^\theta=\Gtheta(\eta;h)\xi$ denote the model value.  Thus $q$ represents
the first Zakharov vector-field component $\eta_t$, not an additional state
variable.  Let $B$ be the batch size, and let hats denote the one-sided real
FFT in the spatial coordinate.
Within the modal band below, $q^\star$ is the projected order-six reference
value for every source.  The accepted C27 model uses the mean per-sample
relative one-sided spectral $L^2$ loss
\begin{equation}
  \ell_{2,b}
  =\frac{\left(\sum_{k=0}^{N/2}
  |\widehat{\widetilde q}^{\,\theta}_{b,k}
   -\widehat{\widetilde q}^{\,\star}_{b,k}|^2\right)^{1/2}}
  {\left[\max\left\{
  \sum_{k=0}^{N/2}|\widehat{\widetilde q}^{\,\star}_{b,k}|^2,
  10^{-12}\right\}\right]^{1/2}},
  \qquad
  \mathcal L_2=\frac1B\sum_{b=1}^B\ell_{2,b}.
  \label{eq:data-loss}
\end{equation}
Here $\widetilde q=q/s_q$ uses the fixed global target scale supplied to the
network.  Away from the numerical floor, $s_q$ and the fixed FFT normalization
cancel.  Equation~\eqref{eq:data-loss} is the implemented one-sided real-FFT
norm: the interior modes are not multiplied by their Parseval multiplicity.
This convention does not affect the argument below because all positive
interior modes carry the same omitted factor.

The single normalization in \eqref{eq:data-loss} allocates accuracy according
to absolute target power.  To make this precise away from the numerical floor,
let
\begin{equation}
  Q_b=\sum_j|\widehat q_{b,j}^\star|^2,
  \qquad
  \mathcal A_b=\{k:\widehat q_{b,k}^\star\ne0\},
  \qquad
  \mathcal Z_b=\{k:\widehat q_{b,k}^\star=0\}.
\end{equation}
For $e_{b,k}=\widehat q_{b,k}^\theta-\widehat q_{b,k}^\star$,
\begin{equation}
  \ell_{2,b}^2
  =\sum_{k\in\mathcal A_b}
  \underbrace{\frac{|\widehat q_{b,k}^\star|^2}{Q_b}}
  _{\displaystyle \beta_{b,k}}
  \left|\frac{e_{b,k}}{\widehat q_{b,k}^\star}\right|^2
  +\frac1{Q_b}\sum_{k\in\mathcal Z_b}|e_{b,k}|^2.
  \label{eq:l2-modal-weighting}
\end{equation}
Thus a nonzero mode containing $10^{-4}$ of the carrier target power
contributes only $10^{-4}$ as much when the two modes have the same fractional
error.  This is
appropriate for total one-step error, but it can hide a weak harmonic or
Benjamin--Feir sideband whose phase-rate bias accumulates over a long rollout.
The second term in \eqref{eq:l2-modal-weighting} penalizes a coefficient that
the model creates where the stored target is exactly zero.

\subsection{Active-mode phase-space-normalized complex loss}
\label{sec:mode-loss}

The quantity called the \emph{mode-balanced loss} in the implementation is
more precisely an active-mode relative error in the modal vector field
$q=\eta_t$.  It does not compare Fourier magnitudes alone.  It compares the
full complex coefficients, so it retains both radial and angular components
of the instantaneous modal motion.

For $1\leq k\leq K$, with $K=128$, define
\begin{equation}
  g_{0,b,k}=k\tanh(\widetilde h_b k),
  \qquad
  \omega_{b,k}^2=g g_{0,b,k},
  \qquad
  \widetilde h_b=\min\{h_b,5\}.
  \label{eq:mode-dispersion}
\end{equation}
The depth clipping agrees with the model. For $h\geq5$, its largest discrepancy
from the infinite-depth symbol is $1-\tanh(5)\simeq9.1\times10^{-5}$ relative
to $k$, attained at $k=1$. Hats in this subsection denote the
forward-normalized real FFT of the physical, denormalized fields. For every
training example, set
\begin{equation}
  P_{b,k}
  =|\widehat q_{b,k}^\star|^2
   +\omega_{b,k}^2|\widehat\eta_{b,k}|^2,
  \qquad
  S_b=\max_{1\leq j\leq K}P_{b,j}.
  \label{eq:mode-phase-space-scale}
\end{equation}
This scale has a direct origin. Suppressing the sample index, the flat-surface
dynamics of mode $k$ at the clipped depth are
\begin{equation}
  \partial_t\widehat\eta_k=\widehat q_k,
  \qquad
  \partial_t\widehat q_k=-\omega_k^2\widehat\eta_k,
  \qquad
  \frac{\dd}{\dd t}
  \left(|\widehat q_k|^2+\omega_k^2|\widehat\eta_k|^2\right)=0.
  \label{eq:linear-modal-oscillator}
\end{equation}
Moreover, if
\begin{equation}
  H_k^{\mathrm{lin}}
  =\frac12\left(g|\widehat\eta_k|^2
  +g_{0,k}|\widehat\xi_k|^2\right)
\end{equation}
denotes the linear Hamiltonian contribution of that mode at the clipped depth,
then
\begin{equation}
  |\widehat q_k|^2+\omega_k^2|\widehat\eta_k|^2
  =2g_{0,k} H_k^{\mathrm{lin}},
  \qquad \widehat q_k=g_{0,k}\widehat\xi_k.
  \label{eq:mode-scale-hamiltonian}
\end{equation}
Consequently, in the linear limit, $P_{b,k}$ is the squared phase-space radius
of the corresponding oscillator. In nonlinear data,
\eqref{eq:mode-phase-space-scale} evaluates that motivated radius formula using
the nonlinear target $q^\star$; it is only a reference scale and is not asserted
to be a separately conserved nonlinear modal energy. The
$\eta$ contribution also prevents a singular relative error when
$\widehat q_k^\star$ passes through a turning point while the oscillator
remains active.

The Hamiltonian relation therefore fixes the units of the denominator only.
This is not an energy-matching or energy-conservation penalty: the numerator
is the full complex vector-field error
$\widehat q^\theta_k-\widehat q^\star_k$, and the reference scale $P_{b,k}$ is
held fixed during optimization. In particular, coefficients with equal
magnitude but different phase are not identified.

The implemented soft activity, relative defect, and robust penalty are
\begin{align}
  a_{b,k}
  &=\frac{P_{b,k}}{P_{b,k}+\alpha S_b+\epsilon_a},
  &\alpha&=10^{-4},\label{eq:mode-activity}\\
  r_{b,k}
  &=\frac{|\widehat q_{b,k}^\theta-\widehat q_{b,k}^\star|^2}
  {P_{b,k}+\beta S_b+\epsilon_a},
  &\beta&=10^{-6},\label{eq:mode-ratio}\\
  \rho(r)
  &=2\left(\sqrt{1+\min\{r,100\}}-1\right),
  &\epsilon_a&=10^{-24}.
  \label{eq:mode-robust}
\end{align}
All quantities derived from the reference state are held fixed under
differentiation.  The loss is
\begin{equation}
  \ell_{\mathrm{mode},b}
  =\frac{\sum_{k=1}^{K}a_{b,k}\rho(r_{b,k})}
  {\max\{\sum_{k=1}^{K}a_{b,k},\epsilon_a\}},
  \qquad
  \mathcal L_{\mathrm{mode}}
  =\frac1B\sum_{b=1}^B\ell_{\mathrm{mode},b}.
  \label{eq:mode-loss}
\end{equation}
For $r\ll1$, $\rho(r)=r+O(r^2)$.  For $P_{b,k}\gg\alpha S_b$,
$a_{b,k}\approx1$, so \eqref{eq:mode-loss} is approximately the average of
$|e_{b,k}|^2/P_{b,k}$ over reference-active modes rather than the target-power-
weighted average in \eqref{eq:l2-modal-weighting}.  Activity is soft,
not binary: a mode with $P_{b,k}=10^{-4}S_b$ has weight $1/2$.  Since $P$ is a
squared scale, this corresponds roughly to one percent of the strongest modal
phase-space radius. Modes below that level fade continuously, while the smaller
denominator floor and the capped pseudo-Huber penalty prevent numerical-noise
coefficients or a single outlier from dominating.

The connection to phase accumulation is exact at the level of the vector
field.  Write a nonzero elevation coefficient as
\begin{equation}
  \widehat\eta_k=A_k e^{\mathrm i\phi_k}.
\end{equation}
Because $q=\eta_t$, any candidate value of $\widehat q_k$ induces
\begin{equation}
  \widehat q_k=e^{\mathrm i\phi_k}
  \left(\dot A_k+\mathrm i A_k\dot\phi_k\right).
\end{equation}
At the same fixed training state, the reference and model therefore satisfy
\begin{equation}
  |\widehat q_k^\theta-\widehat q_k^\star|^2
  =(\delta\dot A_k)^2+A_k^2(\delta\dot\phi_k)^2,
  \label{eq:mode-rate-decomposition}
\end{equation}
where
$\delta\dot A_k=\dot A_k^\theta-\dot A_k^\star$ and
$\delta\dot\phi_k=\dot\phi_k^\theta-\dot\phi_k^\star$.
Thus the complex loss measures errors in both modal amplitude-growth rate and
modal phase rate without estimating or unwrapping a phase angle.  For a
progressive linear mode, $|\widehat q_k^\star|=\omega_k A_k$, and hence
\begin{equation}
  r_k\approx\frac12\left[
  \left(\frac{\delta\dot A_k}{\omega_k A_k}\right)^2
  +\left(\frac{\delta\dot\phi_k}{\omega_k}\right)^2
  \right].
  \label{eq:mode-linear-rate-error}
\end{equation}
A traveling-wave speed defect gives
$\delta\dot\phi_k=-k\,\delta c$ and $c_k=\omega_k/k$, so its phase component is
$r_k\approx\tfrac12(\delta c/c_k)^2$.  A persistent bias accumulates as
$\delta\phi_k(T)=\int_0^T\delta\dot\phi_k(t)\dd t$ and therefore as a spatial
translation.  This explains why a small one-step defect can matter in a long
rollout.

The loss does not pin an absolute position.  Simultaneously translating
$\eta$, $q^\star$, and $q^\theta$ multiplies their Fourier coefficients by the
same unit complex number and leaves \eqref{eq:mode-loss} unchanged.  It instead
penalizes the erroneous modal velocity that would subsequently create a
relative translation.  It is also complementary to the localized tangent
loss below: the tangent loss retains one source-restricted physical-space
projection of $q^\theta-q^\star$, whereas \eqref{eq:mode-loss} is universal and
retains the complete complex error of every reference-active mode.

\subsection{Hadamard identity}

We suppress the fixed depth argument $h$ in this subsection.
Define the vertical and horizontal surface velocities associated with a DNO
evaluation by
\begin{equation}
  B_\theta=
  \frac{\Gtheta(\eta)\xi+\eta_x\xi_x}{1+\eta_x^2},
  \qquad
  V_\theta=\xi_x-B_\theta\eta_x.
  \label{eq:BV}
\end{equation}
For an exact DNO and any surface perturbation $\zeta$, the Hadamard
shape-derivative formula is \cite[Theorem~3.20]{lannes2005}
\begin{equation}
  \D_\eta\G(\eta)[\zeta]\xi
  =-\G(\eta)(\zeta B)-\partial_x(\zeta V).
  \label{eq:hadamard}
\end{equation}
This is precisely the identity that makes the geometric expression for
$\xi_t$ equal to the negative $\eta$-variation of the kinetic energy
\begin{equation}
  K_\theta(\eta,\xi)=\frac12\ip{\xi}{\Gtheta(\eta)\xi}.
  \label{eq:learned-kinetic}
\end{equation}
Hence it is stronger and more local than merely penalizing the drift of the
scalar Hamiltonian along one approximate time step.

We approximate the derivative by a finite secant,
\begin{equation}
  S_\theta=
  \frac{\Gtheta(\eta+\epsilon\zeta)\xi-
        \Gtheta(\eta)\xi}{\epsilon}
  +\Gtheta(\eta)(\zeta B_\theta)
  +\partial_x(\zeta V_\theta).
  \label{eq:hadamard-residual}
\end{equation}
The randomized direction is zero mean and supported on
$0<|k|\leq128$. Its Fourier coefficients are divided by the $H^1$ weight
before RMS normalization, preventing the probe energy from concentrating at
the largest modes. Its physical RMS is matched to
$\max(\mathrm{rms}(\eta-\bar\eta),10^{-3})$, and
$\epsilon$ is sampled log-uniformly from $[10^{-3},3\times10^{-3}]$.

Let $W(k)^2=1+k^2$ and $\project$ denote projection onto the resolved modes.
The normalized consistency loss is
\begin{equation}
  \mathcal L_{\mathrm{shape}}(\theta)
  =\mathbb E_{(\eta,\xi,h),\zeta}
  \frac{\norm{W\project S_\theta}_2^2}
  {\norm{W\project\left[
      \Gtheta(\eta)(\zeta B_\theta)+
      \partial_x(\zeta V_\theta)\right]}_2^2+10^{-12}}.
  \label{eq:shape-loss}
\end{equation}
The finite secant avoids the mixed input--parameter higher-order automatic
differentiation that would arise when optimizing a Jacobian-vector-product
penalty directly.

\subsection{Localized translation-tangent loss}
\label{sec:tangent-loss}

A translated solitary wave has the same Hamiltonian as the original profile,
so scalar energy control cannot determine its propagation speed.  A global
projection onto $\eta_x$ can cancel between counter-propagating or interacting
crests.  C27 therefore uses a local least-squares projection on the three
Tanaka training sources (historical source identifiers 5, 6, and 14).  First
remove the spatial mean from $q^\theta$ and $q^\star$ and set
$e=(q^\theta-\overline q^\theta)-(q^\star-\overline q^\star)$. Let
$\widetilde h=\min\{h,5\}$, and let $K_{\widetilde h}$ denote periodic Gaussian
smoothing with physical width $\widetilde h$. Define
\begin{align}
  C_{\widetilde h}&=K_{\widetilde h}*(\eta_x e),
  &E_{\widetilde h}&=\max\{K_{\widetilde h}*(\eta_x^2),0\},\\
  \delta c_{\widetilde h}(x)
  &=-\frac{C_{\widetilde h}(x)}
  {E_{\widetilde h}(x)+10^{-3}\max_y E_{\widetilde h}(y)+10^{-30}}.
  \label{eq:localized-speed-error}
\end{align}
For the eligible batch indices
$\mathcal I=\{b:\operatorname{source}_b\in\{5,6,14\},
\max_x E_{\widetilde h_b}(x)>10^{-20}\}$, the dimensionless loss is
\begin{equation}
  \ell_{\mathrm{tan,loc},b}
  =\frac{\int E_{\widetilde h_b}|\delta c_{\widetilde h_b}|^2/
  (g\widetilde h_b)\dd x}{\int E_{\widetilde h_b}\dd x},
  \qquad
  \mathcal L_{\mathrm{tan,loc}}
  =\frac{\sum_{b\in\mathcal I}\ell_{\mathrm{tan,loc},b}}
  {\max\{|\mathcal I|,1\}}.
  \label{eq:localized-tangent-loss}
\end{equation}
The $\sqrt{g\widetilde h}$ speed scale makes the local error dimensionless without an
inverse fitted-speed singularity.  Unlike \eqref{eq:mode-loss}, this term
discards the component of $e$ orthogonal to the locally smoothed translation
direction and is inactive outside the designated solitary-wave sources.

\subsection{Combined objective and optimization}

Let $\chi_{16}(s)=1$ when the zero-based optimizer step satisfies
$s\equiv0\pmod{16}$ and let it be zero otherwise. The C27 per-step training
objective is
\begin{equation}
  \boxed{\begin{aligned}
  \mathcal L_{\mathrm{C27}}(\theta;s)
  &={\mathcal L}_2
  +6w(s)\mathcal L_{\mathrm{mode}}
  +10\mathcal L_{\mathrm{tan,loc}}
  +10^{-2}w(s)\chi_{16}(s)\mathcal L_{\mathrm{shape}},\\
  w(s)&=\min\left(\frac{s}{500},1\right).
  \end{aligned}}
  \label{eq:combined-loss}
\end{equation}
The active-mode term is evaluated on every batch and all wave families.  The
tangent term is evaluated on every batch but averages only the eligible
examples.  The shape term is evaluated in float64 on a random global
microbatch of eight examples every 16 optimizer steps. This
selective evaluation makes the shape regularizer affordable while sampling
many independent states and directions over an epoch.

The coefficient 6 in \eqref{eq:combined-loss} is empirical: the ablation below
tests 6 against 0, not its optimality.  The three target-dependent terms use
the same reference values in different geometries.  The mathematical distinction is
that $\mathcal L_2$ controls total DNO error,
$\mathcal L_{\mathrm{mode}}$ approximately equalizes the influence of the
complete complex error over reference-active modal phase-space scales, and
$\mathcal L_{\mathrm{tan,loc}}$ selects a local
translation component on particular sources.  The Hadamard term is the only
one that adds an independent operator-derivative identity.  Thus
\eqref{eq:combined-loss} has four algebraic summands: three target-dependent
losses and one intermittently evaluated identity penalty.  The three-term
description used below refers only to the target-dependent part; it does not
silently absorb or omit $\mathcal L_{\mathrm{shape}}$.

Checkpoint selection uses
\begin{equation}
  \mathcal L_{\mathrm{val}}
  =\mathcal L_2+6\mathcal L_{\mathrm{mode}}
   +10\mathcal L_{\mathrm{tan,loc}},
  \label{eq:c27-validation-loss}
\end{equation}
without a Hadamard secant. The primary C27/C28 deletion comparison uses the
final epoch-40 checkpoint in both arms; epoch 40 also minimizes each arm's own
validation objective.

We optimize for 40 epochs with AdamW, global batch size 1024, learning rate
$2\times10^{-5}$, weight decay $10^{-4}$, a 500-step learning-rate warmup, and
cosine decay. Training is data-parallel over two devices. Because the periodic
shape regularizer depends on the exact optimizer counter, checkpoint
restoration verifies equality of the full train state across replicas,
equality of Adam and schedule counters, and exact agreement between expected
and observed regularizer activations. Checkpoint data are made durable before
atomic publication of metadata.
